\chapter{Literature Review and Related Work}
\label{chap:relatedworks}

\section{Competitor Analysis}
\label{section:competitor-analysis}

\begin{table}[h!]
\centering
\small
\resizebox{\textwidth}{!}{
\begin{tabular}{|p{3cm}|p{3.6cm}|p{3cm}|p{3cm}|p{3cm}|}
\hline
\textbf{Criteria} & \textbf{Proposed System} & \textbf{Somerville et al. (2024)} & \textbf{Jin et al. (2024)} & \textbf{Doroftei et al. (2024)} \\
\hline
\textbf{Primary Focus} & Adaptive UAV training and mission planning & Simulator-based pilot pre-training & VR-based drone learning & Pilot performance evaluation \\
\hline
\textbf{Mission Planning} & QGroundControl integrated & Not supported & Task-based only & Partial support \\
\hline
\textbf{Simulation} & ArduPilot SITL with Unity & Basic simulator training & VR environment & Virtual test environment \\
\hline
\textbf{Gamification} & Gamified UI with leaderboard & Limited & Strong VR engagement & Minimal \\
\hline
\textbf{AI / ML} & K-Means with Random Forest & None & None & AI-assisted evaluation \\
\hline
\textbf{Adaptation} & Adaptive terrain generation & Fixed scenarios & Guided tasks & Static environment \\
\hline
\textbf{Evaluation} & Difficulty-normalized leaderboard & No leaderboard & No leaderboard & Metrics only \\
\hline
\end{tabular}
}
\caption{Competitor Analysis of the Proposed UAV Training Simulator}
\end{table}

In order to understand the current research landscape and identify systems with related capabilities, a competitor analysis is conducted based on recent studies in drone training, virtual learning, and pilot-performance evaluation. Since no single paper was found that integrates all elements of the proposed system, the comparison is based on the closest research works that address important subproblems of the project. These include simulator-based drone pilot training, immersive VR drone teaching, and quantitative pilot performance assessment.

\textbf{Somerville et al.} \cite{Somerville2024} proposed a simulator-based pre-training approach for drone pilots and evaluated whether simulation improves real-world piloting performance. Their study reported that simulator training improved performance, including a 32 percents reduction in mean final displacement, showing that simulation is effective for early drone skill development. However, the work focuses on pre-training effectiveness only and does not include mission-planning interaction, adaptive environments, machine-learning-based skill classification, or competitive gamification features.

\textbf{Jin et al.} \cite{Jin2024} developed the DVRT system, a virtual reality drone teaching platform for beginners in drone operation and programming. The system emphasizes immersive learning, engagement, and practical teaching through interactive VR activities, and the study reports positive learning outcomes for students. This work is relevant to the gamification and learner-engagement aspect of the proposed project. Nevertheless, its emphasis is on VR teaching and programming education rather than UAV mission planning, telemetry-based skill classification, adaptive terrain generation, or fair leaderboard design.

\textbf{Doroftei et al.} \cite{Doroftei2024} introduced a quantitative methodology for assessing drone pilot performance in a realistic virtual test environment. Their framework records both flight metrics and mission-related performance indicators, and the paper also reports a human-performance model, AI co-pilot support, trajectory optimization, and a mission planning tool for pilot assignment. This makes it the closest paper in terms of structured pilot evaluation. However, the study is primarily assessment-oriented and uses standardized environments for measurement rather than adaptive terrain generation for personalized learning. It also does not report a difficulty-normalized leaderboard for fair comparison across different training conditions.

The proposed system differs from these works by integrating the full training loop into one platform. QGroundControl provides the mission-planning interface, where autonomous missions can be created and uploaded to the vehicle, while ArduPilot SITL enables the autopilot to run without hardware in simulation. These components are combined with Unity to create a gamified training interface rather than a purely technical simulation environment.

More importantly, the proposed system adds an AI-based adaptation mechanism that is not clearly present in the selected competitor papers. Flight telemetry and interaction data are used to classify user skill by applying K-Means clustering and Random Forest classification, after which the system generates adaptive terrain suited to the learner's ability level. In addition, the system introduces a difficulty-normalized leaderboard so that users can still be compared fairly even when they train under different terrain and scenario conditions. This makes the proposed system not only a simulator or teaching tool, but also a personalized and competitively balanced UAV training platform. The closest papers identified here cover only separate parts of that idea, which supports the novelty of the proposed system.

\section{Literature Review}
\label{section:literature-review}

Ground Control Systems (GCS) are widely used to support UAV mission planning and monitoring. According to Dardoize et al. [1], a GCS provides essential functionalities such as waypoint-based mission design, telemetry monitoring, and flight data logging. These systems enable users to define flight paths, visualize UAV trajectories, and monitor system status during mission execution. As a result, GCS platforms form the foundation of UAV flight planning and operation. However, despite these capabilities, existing systems are primarily designed for mission execution rather than evaluating the quality of user-defined flight plans or supporting training-oriented feedback. 

Recent studies have also explored the integration of artificial intelligence techniques within training environments to improve learning outcomes. According to Gombolay et al. [2], machine learning can be used to analyze user behavior, predict performance, and support adaptive training. Their work demonstrates that user-generated data can be transformed into meaningful features that reflect decision-making patterns, enabling systems to provide personalized feedback. This concept is particularly relevant to UAV training, where mission planning decisions directly influence flight performance.

In addition to behavior analysis, gamification and immersive technologies have been applied to enhance learning effectiveness. Research on mixed reality and gamified training environments shows that interactive systems can significantly improve user engagement and skill acquisition [5]. These systems create realistic and motivating learning experiences, which are important for training complex tasks such as UAV mission planning. However, many of these approaches focus on general training environments rather than domain-specific applications such as UAV flight planning.

From a UAV system perspective, machine learning has been widely applied to improve system performance. A comprehensive survey on UAV applications highlights the use of machine learning in areas such as path planning, navigation, and decision-making [4]. Furthermore, Baidya et al. [6] explored the use of machine learning techniques in UAV network environments, demonstrating how intelligent algorithms can optimize system behavior under dynamic conditions. While these studies confirm the effectiveness of machine learning in UAV systems, they primarily focus on autonomous system optimization rather than human-centered training.

One of the key techniques for analyzing user behavior is clustering. As discussed by Gombolay et al. [2], K-Means clustering can be used to discover patterns in user interactions without requiring labeled data. This is particularly useful in UAV training, where user skill levels are not explicitly defined. By analyzing features such as trajectory deviation, speed variation, and control stability, clustering can group users into behavior-based categories. These categories can then be used to understand different levels of user performance.

To enable real-time evaluation, supervised learning techniques are required. Breiman [3] introduced the Random Forest algorithm, which is widely used for classification tasks due to its robustness and ability to handle complex data. Random Forest can effectively classify users based on extracted features and provide insights into which factors most influence performance. This makes it suitable for UAV training systems, where interpretability and reliability are important.

In addition, research on artificial intelligence in serious games has shown that machine learning can support decision-making and adaptive learning processes [7]. These systems can adjust training scenarios based on user behavior, improving learning efficiency and maintaining engagement. This reinforces the potential of combining machine learning with simulation-based training systems.

In conclusion, while existing UAV systems provide strong capabilities in mission planning, simulation, and autonomous control, there remains a gap in supporting human-centered training and behavior evaluation. Ground Control Systems enable effective flight planning, and simulation environments offer realistic testing conditions, but they lack intelligent mechanisms to analyze user performance. At the same time, machine learning techniques such as K-Means clustering and Random Forest classification have been proven effective in behavior analysis and performance prediction in training systems. Therefore, integrating these approaches into a unified UAV training platform is both feasible and justified. The proposed system builds upon these established foundations by combining flight planning, simulation, and machine learning to analyze user behavior and provide adaptive feedback, ensuring a practical and achievable solution for improving UAV training outcomes.
